\documentclass[10pt,a4paper]{amsart}
\usepackage[margin=19mm]{geometry}
\usepackage{amsmath,amssymb,mathtools}
\usepackage[scr=rsfs,cal=euler]{mathalfa}
\usepackage{xfrac}
\usepackage[unicode,hidelinks,bookmarks=false]{hyperref}
\newcommand{\ie}{i.e.\ }
\newcommand{\cf}{cf.\ }
\newcommand{\resp}{respectively\ }
\newcommand{\Subsection}{\subsection}
\providecommand{\nobreakdash}{\nobreak\hbox{-}}
\setlength{\emergencystretch}{2em}
\multlinegap0pt
\usepackage{mathtools}
\usepackage{amssymb}
\newtheorem{construction}{Construction}
\newtheorem{proposition}{Proposition}
\newcommand{\C}{\mathbb{C}}
\DeclareMathOperator{\Cl}{Cl}
\newcommand{\Q}{\mathbb{Q}}
\DeclareMathOperator{\Stab}{Stab}
\newcommand{\Z}{\mathbb{Z}}
\DeclareMathOperator{\im}{im}
\title{A canonicity criterion for toric varieties \\ and the classification of canonical 4-simplices}
\author{Marco Ghirlanda}
\date{Source: arXiv:2603.21198v1 (CC BY 4.0)}
\begin{document}
\maketitle

\noindent\emph{Source and licence.} A canonicity criterion for toric varieties \\ and the classification of canonical 4-simplices. Version arXiv:2603.21198v1 (CC BY 4.0), distributed by arXiv under the Creative Commons Attribution 4.0 International licence, http://creativecommons.org/licenses/by/4.0/, as recorded in the arXiv licence metadata for this exact version and shown on the abstract page https://arxiv.org/abs/2603.21198v1. Full licence notice: \texttt{licenses/arxiv-2603.21198-CC-BY-4.0.txt}.\par\smallskip

\noindent\textit{Editorial scope.} Counted unit: the article's proposition prop:eigenvals (source0.tex lines 394--410). The article's complete original proof of it is delivered with it (source0.tex lines 413--468, one \texttt{proof} environment of 56 lines). No mathematics is written, repaired or re-derived: the statement and the proof are the article's own lines, in the article's own notation, and no retained line is substituted or re-flowed.  Retained uncounted as the source-local context the proof needs: lines 267--302; lines 354--377.  Nothing was omitted from any retained span, so the package carries no omission record.  No retained line contains a citation, so the wrapper carries no bibliography and no bibliography entry.  No retained line contains a figure environment, an image-inclusion command or drawing code, so no image is delivered and the package declares no asset dependency.  Editorial formatting only, in addition: an AMS \texttt{amsart} preamble that restates the article's own declarations and macros used by the retained lines, namely the environments \texttt{construction}, \texttt{proposition} on separate counters, together with the standard packages those lines need.  The retained environments print in this document as Construction 1, Proposition 1 in order of appearance, whereas the article prints the same items with the numbering of its own full document.  The complete native source of the article as distributed in the arXiv e-print for this exact version is bundled unchanged as \texttt{sources/196924/source0.tex}.
\begin{construction}\label{constr:slice}
Let $\Sigma$ be a simplicial fan in~$\Z^n$ with generator
matrix $P$.
Consider the associated toric variety $X$ with Cox's quotient
presentation $p \colon \hat X \to X$.
For any $n$-dimensional $\sigma \in \Sigma$,
the \emph{affine slice} at the distinguished point
$\hat x_{\hat \sigma} \in \hat X$ is
$$
Y_\sigma
\coloneqq
\{ x \in \C^r; \ x_i=1 \text{ for all } i = 1, \ldots, r \text{ with } \  v_i \not\in \sigma \}
\subseteq
\hat X.
$$
Note that $Y_\sigma$ is an $n$-dimensional
affine subspace of $\C^r$.
For the characteristic quasitorus $H$ and the
\emph{stabilizer} $H_\sigma := \Stab_H(Y_\sigma) \subseteq H$
of $Y_\sigma \subseteq \hat X$, we have
$$
H \cdot Y_\sigma = \hat X(\hat\sigma),
\qquad
X(\sigma) = \hat X(\hat\sigma)/H = Y_\sigma/H_\sigma.
$$
Moreover, the stabilizer $H_\sigma$ of $Y_\sigma$
equals the isotropy group $H_{\hat x_{\hat \sigma}}$
and thus its character group is given by
$$
\mathbb{X}(H_\sigma) = \Cl(X,x_\sigma),
$$
where $\Cl(X,x_\sigma)$ is the local class group of the distinguished
point $x_\sigma \in X$, that means the group of Weil divisors
modulo those being principal near $x_\sigma$. In particular,
as $X$ is $\Q$-factorial, $H_\sigma$ is finite.
\end{construction}

\begin{construction}\label{constr:ev1}
Consider a simplicial fan $\Sigma$ in $\Z^n$ with generator
matrix $P$ and associated degree matrix $Q$.
Let $\sigma \in \Sigma$ be an $n$-dimensional cone
and let $1 \le i_1 < \ldots < i_k \le r$ be the
indices with $v_{i_j} \not \in \sigma$. These give
us matrices:
$$
W_\sigma \ := \ \begin{bmatrix} w_{i_1} \ldots w_{i_k} \end{bmatrix}^T,
\qquad
E_\sigma \ := \ \begin{bmatrix} \eta_{i_1} \ldots \eta_{i_k} \end{bmatrix}^T.
$$
where $\omega_i = (w_i,\eta_i) \in K = \Z^k \times \Gamma$ is the $i$-th column
of $Q$. For any $\eta \in \Gamma$, we write $\eta = (e_1, \ldots, e_s)$
with $0 \le e_l <  \mu_l$ and set $b(\eta) = (e_1/\mu_1,\ldots,e_s/\mu_s)$.
Then we obtain a monomorphism
$$
\begin{array}{rcl}
\Phi \colon \Z^k/\mathrm{im}(W_\sigma) \times \Gamma & \to & \Q^k/\Z^k \times \Q^s/\Z^s,
\\[5pt]                                                     
(c,\eta) & \mapsto & \left(W_\sigma^{-1} \cdot (c-E_\sigma \cdot b(\eta)), \, b(\eta) \right).
\end{array}
$$
\end{construction}

\begin{proposition}\label{prop:eigenvals}
Consider the setting of Construction~\ref{constr:ev1}.   
With the monomorphism $\Phi$ we obtain an isomorphism of groups
$$
\begin{array}{rcl}
\mathrm{im}(\Phi) & \to & H_\sigma,
\\[5pt]                                                     
(a,b) & \mapsto & h(a,b) := (\exp(2 \pi i a_1), \ldots, \exp(2 \pi i a_k), \, \exp(2 \pi i b_1), \ldots, \exp(2 \pi i b_s)).
\end{array}
$$
Denote by $\alpha_j(a,b)$ the fractional parts of $\omega_{m_j} \cdot (a,b)$,
where $1 \le m_1 < \ldots < m_n \le r$ are the indices with $v_{m_j} \in \sigma$.
Then the eigenvalues of $h(a,b) \in H_\sigma$ acting on~$Y_\sigma \cong \C^n$ are
$$
\exp(2 \pi i \alpha_1(a,b)), \ldots, \exp(2 \pi i \alpha_n(a,b)).
$$
\end{proposition}

\begin{proof}
We begin with a short preparation.
Write the transposed degree matrix as $Q^T= [W,E]$, where $W$
(resp. $E$) has the upper $k$ (resp. lower $s$) rows of $Q$ as
their columns. Denote by $\Q_\mu^s=\Q_{\mu_1}\times\dots\times\Q_{\mu_s}\subset\Q^s$ the subgroup of all tuples $(b_1/\mu_1,\dots,b_s/\mu_s)$ with $b_i\in\Z$. Consider now the homomorphism 
$$
\Psi \colon
\Q^k/\Z^k \times \Q_\mu^s/\Z^s \ \to \ \Q^r/\Z^r,
\qquad
(a,b) \ \mapsto \ Q^T \cdot (a, b) =  W \cdot a + E \cdot b.
$$  
Then $\Psi$ tells us how the elements $h = (t,\zeta)$ of
finite order in $H$ act on $\C^r$.
Any such element can be written uniquely as 
$$
h
\ = \
(t,\zeta)
\ = \
\left(
\exp(2 \pi i a_1), \ldots, \exp(2 \pi i a_k),
\exp(2 \pi i b_1), \ldots, \exp(2 \pi i b_s)
\right),
$$
where $(a_1,\ldots, a_k) \in \Q^k/\Z^k$ and $(b_1,\ldots, b_s) \in \Q_\mu^s/\Z^s$.
With the rows $(w_l,\eta_l)$ of the transpose $Q^T$ and the class of
$\mathfrak{a}_l := (w_l,\eta_l) \cdot (a,b)$ in $\Q/\Z$ we have
$$
h \cdot z
\ = \
(t^{w_1}\zeta^{\eta_1}z_1, \ldots, t^{w_r}\zeta^{\eta_r}z_r)
\ = \
\left(\exp(2 \pi i\mathfrak{a}_1)z_1, \ldots, \exp(2 \pi i \mathfrak{a}_r)z_r \right).
$$

Now, by Construction~\ref{constr:slice}, the group $H_\sigma \subseteq H$
consists precisely of the finite order elements of $H$ having eigenvalue
$1$ at the coordinates $l$ with $v_l \not \in \sigma$.
Using the definitions of $\Phi$ and $\Psi$, we see that any $h(a,b)$ with $(a,b)\in\im(\Phi)$
is actually an element of~$H_\sigma$.
Obviously, $\Phi$ is injective.
For surjectivity, let $h = (t,\zeta) \in H_\sigma$ be given.
Write $t_l=\exp(2\pi i a_l)$ and  $\zeta_l=\exp(2\pi i b_l)$
with $a_l \in \Q/Z$ and $b_l \in \Q_{\mu_l}/\Z$.
We claim
$$
c := W_\sigma \cdot a + E_\sigma \cdot b \in \Z^k.
$$
Indeed, as $h$ has eigenvalue 1 at the coordinates $l$
with $v_l \not \in \sigma$, the corresponding exponents
$\mathfrak{a}_l = (w_l,\eta_l) \cdot (a,b)$ must be integral.
Thus,  $c_l=\mathfrak{a}_l$ gives the claim.
Finally, using again the above presentation of $h \cdot z$ for
$z \in \C^r$, we see that $H_\sigma$ acts on the affine slice
$Y_\sigma$ as claimed in the assertion, with $\alpha_j=\mathfrak{a}_{m_j}$.
\end{proof}

\end{document}
